% Generated deterministically from the audited Markdown source.
% Responsible author: Carptopus.
\documentclass[11pt]{article}
\usepackage[a4paper,margin=27mm]{geometry}
\usepackage[T1]{fontenc}
\usepackage[utf8]{inputenc}
\usepackage{lmodern}
\usepackage{microtype}
\usepackage{amsmath,amssymb,amsthm,mathtools}
\usepackage{needspace}
\usepackage{xcolor}
\usepackage[colorlinks=true,linkcolor=blue!55!black,citecolor=blue!55!black,urlcolor=blue!65!black]{hyperref}
\usepackage{xurl}
\usepackage{fancyhdr}

\newtheorem{theorem}{Theorem}[section]
\newtheorem{lemma}[theorem]{Lemma}
\newtheorem{corollary}[theorem]{Corollary}

\pagestyle{fancy}
\fancyhf{}
\fancyhead[L]{Carptopus}
\fancyhead[R]{Reduced composition polynomials}
\fancyfoot[C]{\thepage}
\setlength{\headheight}{14pt}
\setlength{\parindent}{0pt}
\setlength{\parskip}{0.44em}
\setlength{\emergencystretch}{4em}
\urlstyle{same}

\title{Log-concavity of reduced composition polynomials\\
via adjacent B-spline refinement columns}
\author{Carptopus\\\href{mailto:carptopus@163.com}{carptopus@163.com}}
\date{4 September 2026}

\hypersetup{
  pdftitle={Log-concavity of reduced composition polynomials via adjacent B-spline refinement columns},
  pdfauthor={Carptopus},
  pdfsubject={Log-concavity and unimodality of Ardila-Doker reduced composition polynomials},
  pdfkeywords={composition polynomial, log-concavity, unimodality, discrete B-spline, knot refinement, total positivity}
}

\begin{document}
\maketitle

\begin{center}
\fcolorbox{black!35}{black!4}{\parbox{0.88\textwidth}{\centering\small
Internal candidate manuscript, version 0.1-beta. The mathematical proof and complete
Markdown manuscript have passed project-internal zero-trust audits. External mathematical
review and formal peer review are pending.}}
\end{center}

\begin{abstract}


Ardila and Doker associated a reduced composition polynomial $f_c(q)$ with every positive integer
composition $c$. They proved positivity of its coefficients and asked whether those coefficients are
unimodal or, more strongly, log-concave. We prove the stronger assertion for every composition.
The first ingredient identifies the coefficient vector, up to a positive scalar, with a column of a
B-spline knot-refinement matrix. Total nonnegativity, constant reproduction, and variation
diminution then give unimodality. For log-concavity, deleting the first or last endpoint of the
composition produces two adjacent columns of a lower-degree refinement matrix. Their adjacent
$2\times2$ minors are nonnegative. An exact identity converts those minors into the log-concavity
defects of $f_c$, with the remaining sign case closed by unimodality. Thus both the unimodality
question and the stronger log-concavity question have affirmative answers.

\end{abstract}

\section{Introduction}


Let $c=(d_1,\ldots,d_r)$ be a composition of $n$ into positive parts. Ardila and Doker introduced a
composition polynomial $g_c(q)$ in their study of subdivisions of lifted generalized permutahedra
\cite{ardila-doker}. They showed that

\[
g_c(q)=(1-q)^r f_c(q),
\]

where the reduced composition polynomial $f_c(q)$ has positive coefficients. Their Question 7.5
asks whether these coefficients are unimodal; immediately afterwards they ask whether the stronger
property of log-concavity always holds. The original paper verified the expected shape in a large
finite family but left both questions open.

Novelli and Thibon later gave a combinatorial interpretation of $n!f_c(q)$ as a generating
polynomial for a statistic on permutitions \cite{novelli-thibon}. That interpretation explains coefficient
integrality and positivity for the scaled polynomial $n!f_c(q)$, but it does not by itself provide
the required second-order coefficient inequalities.

Our proof uses a different interface. The explicit coefficient formula for $f_c$ is a
rising-factorial divided difference. This is exactly the formula for a discrete B-spline, or
equivalently for one column of a B-spline knot-refinement matrix \cite{cohen-lyche-riesenfeld}. The one-column consequence is
unimodality. The stronger result comes from looking at two columns at once: the reduced polynomials
obtained by deleting the final and initial parts of $c$ form adjacent refinement columns, so total
nonnegativity controls their adjacent determinants. A short identity turns this determinant
inequality into log-concavity of the original coefficient sequence.

\Needspace{0.18\textheight}
\begin{theorem}
\label{thm:main}

For every positive integer composition $c$, the coefficient sequence of the reduced composition
polynomial $f_c(q)$ is log-concave. In particular, it is unimodal.

\end{theorem}

Here a finite nonnegative sequence $(a_i)$ is log-concave when

\[
a_i^2\geq a_{i-1}a_{i+1}
\]

at every internal index. All coefficient sequences below are extended by zero outside their natural
support.

The spline results used here are classical. The contribution is the exact identification of the
composition-polynomial coefficients with refinement columns and the adjacent-column determinant
argument that proves Theorem~\ref{thm:main}. In particular, the log-concavity of a continuous B-spline as a
function on its support \cite{floater} is distinct from the discrete coefficient log-concavity proved here.

\section{Reduced composition polynomials and discrete B-splines}


Let

\[
c=(d_1,\ldots,d_r),\qquad n=d_1+\cdots+d_r,
\]

and write its partial sums as

\[
0=\beta_0<\beta_1<\cdots<\beta_r=n.
\]

Write

\[
f_c(q)=\sum_{i=0}^{n-r}F_iq^i.
\]

Ardila and Doker's coefficient formula is

\[
F_i=
\sum_{\substack{0\leq j\leq r\\ \beta_j\leq i}}
\frac{(-1)^j}{[\beta_j]}
\binom{r+i-\beta_j-1}{r-1},
\qquad
[\beta_j]=\prod_{\ell\ne j}|\beta_\ell-\beta_j|.
\tag{2.1}
\]

We now place (2.1) in the discrete-spline setting. For strictly increasing integer knots
$z_0<\cdots<z_s$, define the unnormalized discrete B-spline of degree $s-1$ and unit mesh by

\[
U^{s-1}(x;z_0,\ldots,z_s)
=(-1)^s
\bigl[(x-\,\cdot\,)^{\overline{s-1}}\mathbf 1_{x>\,\cdot\,}
\bigr][z_0,\ldots,z_s],
\tag{2.2}
\]

where $y^{\overline m}=y(y+1)\cdots(y+m-1)$ and the final brackets denote a divided difference in
the knot variable. Its normalized version is

\[
V^{s-1}(x;z_0,\ldots,z_s)
=(z_s-z_0)U^{s-1}(x;z_0,\ldots,z_s).
\tag{2.3}
\]

\Needspace{0.18\textheight}
\begin{lemma}
\label{lem:discrete-column}

For $0\leq i\leq n-r$,

\[
V^{r-1}(i+1;\beta_0,\ldots,\beta_r)=n(r-1)!F_i.
\tag{2.4}
\]

The values on the left vanish on the surrounding integer grid outside this index range.

\end{lemma}

\begin{proof}

The divided-difference denominator at $\beta_j$ has sign $(-1)^{r-j}$ and absolute value
$[\beta_j]$. Moreover,

\[
\frac{(i+1-\beta_j)^{\overline{r-1}}}{(r-1)!}
=\binom{r+i-\beta_j-1}{r-1}.
\]

Substituting the knots $\beta_0,\ldots,\beta_r$ into (2.2), and then using (2.3), gives (2.4)
term by term from (2.1).

It remains to verify that the divided difference in (2.4) is a refinement coefficient, rather than
merely a formally similar discrete-spline value. Embed the local knot window in

\[
\tau=(-r,\ldots,-1,\beta_0,\ldots,\beta_r,n+1,\ldots,n+r)
\]

and refine to the unit knot vector

\[
t=(-r,-r+1,\ldots,n+r).
\]

At a fine-grid left endpoint $i$, the refinement formula is the divided difference of a truncated
polynomial kernel. More precisely, Theorem 1 and equations (3.20)--(3.23) of
Cohen--Lyche--Riesenfeld \cite{cohen-lyche-riesenfeld}, specialized to degree $r-1$, unit fine mesh, and the selected coarse
window, give its refinement coefficient as

\[
\alpha_i
=n[\beta_0,\ldots,\beta_r]
\left(
\mathbf1_{y>i}\prod_{h=1}^{r-1}(y-i-h)
\right),
\tag{2.5}
\]

where the divided difference is taken in $y$. Put $x=i+1$ and

\[
R(y)=(x-y)^{\overline{r-1}}.
\]

At every integer knot,

\[
\prod_{h=1}^{r-1}(y-i-h)=(-1)^{r-1}R(y).
\tag{2.6}
\]

The polynomial $R$ has degree $r-1$, so its order-$r$ divided difference on the $r+1$ knots is
zero. On integer knots, $y>i$ and $x>y$ are complementary conditions. When $r>1$, the possible
boundary value $y=x$ vanishes because $R(x)=0$; when $r=1$, strict complementarity holds directly.
Therefore

\[
[\beta_0,\ldots,\beta_r](\mathbf1_{y>i}R)
=-[\beta_0,\ldots,\beta_r](\mathbf1_{x>y}R).
\tag{2.7}
\]

Substitution of (2.6)--(2.7) into (2.5) gives

\[
\alpha_i=V^{r-1}(i+1;\beta_0,\ldots,\beta_r),
\]

including its sign and normalization. Finally, if $i<0$, the truncated kernel in (2.5) agrees at
all coarse knots with one polynomial of degree $r-1$, so its order-$r$ divided difference vanishes.
If $i>n-r$, every potentially nonzero term contains a zero factor. Hence the nonzero core has
precisely the indices in (2.4).
\end{proof}


\section{Knot refinement and unimodality}


We recall only the refinement properties needed below. Suppose a coarse knot vector is refined by
successive knot insertions, without changing the spline degree. If $p$ is the coarse control vector
and $d$ the refined control vector, write

\[
d=Kp.
\]

For one knot insertion, every new control coefficient is either copied from an old coefficient or
is a convex combination of two adjacent old coefficients. In the convention above, the rows of
$K$ are fine indices and its columns are coarse indices. A one-step matrix is therefore a
nonnegative bidiagonal matrix with row sums one. It is totally nonnegative, and products preserve
total nonnegativity and row sums.

For a real vector $v$, let $S^-(v)$ be the number of sign changes after its zero entries are deleted.
We use the standard variation-diminishing theorem for totally nonnegative matrices \cite[Chapter~3]{pinkus}:

\[
S^-(Kv)\leq S^-(v).
\tag{3.1}
\]

\Needspace{0.18\textheight}
\begin{lemma}
\label{lem:unimodal-column}

Every column of a finite B-spline knot-refinement matrix, with zeros retained outside its local
support, is unimodal.

\end{lemma}

\begin{proof}

Let $\alpha=Ke_s$ be a column and choose $0<\lambda<1$. Since $K\mathbf 1=\mathbf 1$,

\[
\alpha-\lambda\mathbf 1=K(e_s-\lambda\mathbf 1).
\tag{3.2}
\]

The vector on the right before applying $K$ has sign pattern

\[
-\;\cdots\;-\;+\;-\;\cdots\;-,
\]

so it has at most two sign changes. Equation (3.1) gives the same upper bound for the left side of
(3.2).

If $\alpha$ were not unimodal, two higher values would be separated by a strictly lower value.
Choose $\lambda$ strictly between the valley and the smaller of the two higher values. Including
the zero entries on the two sides of the local support, $\alpha-\lambda\mathbf1$ would then have a
subpattern

\[
-,+,-,+,-,
\]

and hence at least four sign changes, a contradiction.
\end{proof}


\Needspace{0.18\textheight}
\begin{corollary}
\label{cor:unimodal}

For every positive integer composition $c$, the coefficient sequence of $f_c(q)$ is unimodal.

\end{corollary}

\begin{proof}

By Lemma~\ref{lem:discrete-column}, the coefficient sequence is a positive scalar multiple of one refinement column.
Apply Lemma~\ref{lem:unimodal-column}.
\end{proof}


Thus the first part of Ardila--Doker Question 7.5 follows from a one-column shape property. To reach
log-concavity we use two neighboring columns.

\section{Adjacent endpoint-deletion columns}


Assume for now that $r\geq2$. Define

\[
A(q)=f_{(d_1,\ldots,d_{r-1})}(q),
\qquad
B(q)=q^{d_1}f_{(d_2,\ldots,d_r)}(q),
\tag{4.1}
\]

and write their zero-padded coefficients as $A_i$ and $B_i$. The shift in $B(q)$ places both
sequences on the global unit grid determined by $0,1,\ldots,n$.

The two knot windows

\[
(\beta_0,\ldots,\beta_{r-1})
\quad\text{and}\quad
(\beta_1,\ldots,\beta_r)
\tag{4.2}
\]

are adjacent windows in a common degree-$(r-2)$ coarse B-spline basis. Refine the padded coarse knot
vector

\[
\tau=(-(r-1),\ldots,-1,\beta_0,\ldots,\beta_r,n+1,\ldots,n+r-1)
\]

to the unit grid

\[
t=(-(r-1),-(r-2),\ldots,n+r-1).
\]

Lemma~\ref{lem:discrete-column}, applied to the two shorter compositions, identifies the corresponding adjacent columns
as

\[
\alpha_i^L=\beta_{r-1}(r-2)!A_i,
\qquad
\alpha_i^R=(n-\beta_1)(r-2)!B_i.
\tag{4.3}
\]

Both factors are strictly positive. The refinement matrix is totally nonnegative, and its rows and
columns in (4.3) are in increasing order. Consequently every adjacent $2\times2$ minor gives

\[
M_i:=A_iB_{i+1}-A_{i+1}B_i\geq0.
\tag{4.4}
\]

The orientation in (4.4) is important: rows are consecutive fine-grid indices, and the left window
in (4.2) precedes the right window.

\section{The determinant identity}


Continue to write

\[
F(q)=f_c(q)=\sum_iF_iq^i,
\]

with $F_i=0$ outside the natural support, and set

\[
\Delta_i=F_i-F_{i-1},
\qquad
L_i=F_i^2-F_{i-1}F_{i+1}.
\tag{5.1}
\]

The quantity $L_i$ is the log-concavity defect at $i$.

\Needspace{0.18\textheight}
\begin{lemma}
\label{lem:endpoint}

The endpoint-deletion polynomials satisfy

\[
A(q)-B(q)=n(1-q)F(q)
\tag{5.2}
\]

and

\[
B(q)=q\bigl(rF(q)-(1-q)F'(q)\bigr).
\tag{5.3}
\]

Equivalently, with $k=r-1$,

\[
A_i-B_i=n\Delta_i,
\qquad
B_i=(k+i)F_{i-1}-iF_i.
\tag{5.4}
\]

\end{lemma}

\begin{proof}

The explicit partial-fraction formula of Ardila--Doker is

\[
g_c(q)=\sum_{j=0}^r\frac{(-1)^jq^{\beta_j}}{[\beta_j]}.
\tag{5.5}
\]

Its endpoint-deletion formulas (Corollary 7.3 of \cite{ardila-doker}) read

\[
g_{(d_1,\ldots,d_{r-1})}(q)
=\sum_{j=0}^r
\frac{(-1)^j(n-\beta_j)q^{\beta_j}}{[\beta_j]},
\tag{5.6}
\]

and

\[
q^{d_1}g_{(d_2,\ldots,d_r)}(q)
=-\sum_{j=0}^r
\frac{(-1)^j\beta_jq^{\beta_j}}{[\beta_j]}.
\tag{5.7}
\]

Adding the right side of (5.6) to the negative of the right side of (5.7) gives $ng_c(q)$.
Since $g_c=(1-q)^rF$, while the two truncated polynomials in (5.6)--(5.7) have $r-1$ parts,
division by $(1-q)^{r-1}$ gives (5.2).

Equations (5.7) and (5.5) also give

\[
q^{d_1}g_{(d_2,\ldots,d_r)}(q)=-qg_c'(q).
\]

Differentiating $g_c(q)=(1-q)^rF(q)$ therefore gives

\[
(1-q)^{r-1}B(q)
=q(1-q)^{r-1}\bigl(rF(q)-(1-q)F'(q)\bigr),
\]

which proves (5.3). Extracting coefficients from (5.2) and (5.3) yields (5.4).
\end{proof}


\Needspace{0.18\textheight}
\begin{lemma}
\label{lem:determinant}

For every integer $i$,

\[
M_i=n\bigl(kL_i-\Delta_i\Delta_{i+1}\bigr),
\qquad k=r-1.
\tag{5.8}
\]

\end{lemma}

\begin{proof}

Using $A_i=B_i+n\Delta_i$ in (4.4) gives

\[
M_i=n\bigl(\Delta_iB_{i+1}-\Delta_{i+1}B_i\bigr).
\]

Substitute the second identity in (5.4), once at $i$ and once at $i+1$. Expanding and collecting
the coefficients of $F_{i-1}F_i$, $F_i^2$, and $F_iF_{i+1}$ yields

\[
\Delta_iB_{i+1}-\Delta_{i+1}B_i
=k(F_i^2-F_{i-1}F_{i+1})
-(F_i-F_{i-1})(F_{i+1}-F_i),
\]

which is (5.8).
\end{proof}


\section{Proof of the main theorem}


We first treat $r\geq2$. By Corollary~\ref{cor:unimodal}, $(F_i)$ is a nonnegative unimodal sequence. Hence its
first differences can change sign only from nonnegative to nonpositive; a negative difference can
never be followed later by a positive one.

Fix an internal index $i$. If

\[
\Delta_i\Delta_{i+1}\geq0,
\]

then (4.4) and (5.8) imply

\[
kL_i\geq\Delta_i\Delta_{i+1}\geq0.
\]

Because $k=r-1>0$, this gives $L_i\geq0$.

It remains to consider $\Delta_i\Delta_{i+1}<0$. Unimodality excludes
$\Delta_i<0<\Delta_{i+1}$, so necessarily

\[
\Delta_i>0>\Delta_{i+1}.
\]

Thus $F_i$ is strictly larger than both neighboring coefficients. Since all coefficients are
nonnegative,

\[
F_i^2>F_{i-1}F_{i+1},
\]

and again $L_i\geq0$.

When $r=1$, one has directly

\[
f_{(n)}(q)=\frac{1+q+\cdots+q^{n-1}}{n},
\]

whose coefficients are log-concave. If $r\geq2$ but $n-r\leq1$, the support has at most two terms
and there is no internal inequality to check. With zero padding, the two endpoint inequalities are
also immediate. This proves Theorem~\ref{thm:main} for every positive integer composition.
\hfill$\square$


\section{Examples and scope}


For $n=4$, scaled reduced composition polynomials listed by Novelli--Thibon include

\[
4!f_{(4)}(q)=6+6q+6q^2+6q^3,
\]

\[
4!f_{(3,1)}(q)=2+4q+6q^2,
\]

and

\[
4!f_{(2,2)}(q)=3+6q+3q^2.
\]

The first is flat, the second is increasing, and the third has an interior peak; all three are
log-concave. The proof above treats these patterns uniformly. It also permits equality in the
log-concavity inequalities, so no strictness assertion is intended.

The result concerns the coefficient sequence of $f_c(q)$. It does not assert real-rootedness,
ultra-log-concavity, or a Pólya-frequency property. Nor does it claim that arbitrary columns of an
arbitrary nonnegative or stochastic matrix are log-concave. The determinant argument uses the
specific pair of adjacent endpoint-deletion columns together with the exact identity (5.8).

\section{Computational verification}


The proof is symbolic and does not depend on finite computation. The exact-arithmetic script
\texttt{test\_truncation\_pencil.py} uses (2.1) to construct the coefficients; it does not independently
prove that source formula. For all $8{,}178$ nontrivial compositions of total size at most $13$,
it checked:

\begin{itemize}
\item the coefficient form of (5.2), including the zero-padded endpoint shifts;
\item $53{,}157$ adjacent-column determinants $M_i\geq0$;
\item the same $53{,}157$ instances of identity (5.8);
\item $36{,}813$ internal log-concavity defects.

\end{itemize}

No failure occurred. This finite calculation is an algebraic regression test, not a substitute for
the general proof.

The separate script \texttt{verify\_composition\_unimodality.py} checked positivity, unimodality, and
log-concavity in $75{,}535$ test instances, representing $74{,}913$ distinct compositions. It also
compared the explicit coefficients with actual SciPy knot insertions on $199$ distinct fixed-seed
samples and included deliberately shifted and reversed negative controls. Those checks support the
implementation of Lemma~\ref{lem:discrete-column} but are likewise not used as a logical premise of Theorem~\ref{thm:main}.

\section{Relation to prior work}


Ardila--Doker \cite{ardila-doker} introduced the polynomials, proved coefficient positivity, and posed unimodality
and log-concavity as questions. Their Corollary 9.2 also proves log-concavity of the linear-extension
numbers $N_j$ in a Bernstein-basis expansion of $g_c(q)$. That is a different coefficient sequence
in a different basis from the power-basis coefficients $F_i$ of $f_c(q)$ treated here.
Novelli--Thibon \cite{novelli-thibon} supplied a combinatorial interpretation through permutitions and noncommutative
symmetric functions. Cohen--Lyche--Riesenfeld \cite{cohen-lyche-riesenfeld} developed the discrete B-spline refinement
formulas and subdivision machinery used in Lemma~\ref{lem:discrete-column}. Floater \cite{floater} proved log-concavity of continuous
B-spline functions; that theorem concerns function values on a continuous support rather than the
discrete refinement coefficients considered here. The variation-diminishing implication used in
Section 3 is the standard total-positivity result recorded, for example, by Pinkus \cite{pinkus}.

The one-column refinement argument proves unimodality. The two-column argument then strictly
strengthens it to log-concavity. These are therefore presented as one result rather than as two
independent contributions.

A targeted literature refresh through September 3, 2026 found no publication directly proving
log-concavity for all reduced composition polynomials. This is a bounded novelty search, not a claim
that all unindexed or unpublished work has been excluded.

\section*{AI-assisted research and writing disclosure}


The author used OpenAI Codex as an AI-assisted tool for literature organization, proof exploration,
exact verification, adversarial auditing, and manuscript drafting. Carptopus is responsible for the
content of the manuscript.

\begin{thebibliography}{9}
\footnotesize
\raggedright

\bibitem{ardila-doker}
F.~Ardila and J.~Doker,
\emph{Lifted generalized permutahedra and composition polynomials},
Advances in Applied Mathematics \textbf{50} (2013), 607--633.
\href{https://doi.org/10.1016/j.aam.2013.01.005}{doi:10.1016/j.aam.2013.01.005}.

\bibitem{cohen-lyche-riesenfeld}
E.~Cohen, T.~Lyche, and R.~Riesenfeld,
\emph{Discrete B-splines and subdivision techniques in computer-aided geometric design and computer graphics},
Computer Graphics and Image Processing \textbf{14} (1980), 87--111.
\href{https://doi.org/10.1016/0146-664X(80)90040-4}{doi:10.1016/0146-664X(80)90040-4}.

\bibitem{floater}
M.~S.~Floater,
\emph{Log-concavity of B-splines},
Journal of Approximation Theory \textbf{300} (2024), 106042.
\href{https://doi.org/10.1016/j.jat.2024.106042}{doi:10.1016/j.jat.2024.106042}.

\bibitem{novelli-thibon}
J.-C.~Novelli and J.-Y.~Thibon,
\emph{On composition polynomials},
Journal of Combinatorial Theory, Series A \textbf{152} (2017), 1--9.
\href{https://doi.org/10.1016/j.jcta.2017.06.005}{doi:10.1016/j.jcta.2017.06.005}.

\bibitem{pinkus}
A.~Pinkus,
\emph{Totally Positive Matrices},
Cambridge Tracts in Mathematics 181, Cambridge University Press, 2010.
\href{https://doi.org/10.1017/CBO9780511691713}{doi:10.1017/CBO9780511691713}.

\end{thebibliography}

\end{document}
